\paragraph{二维切片塌缩与原子 Witt 校正的变量分裂}
在二维加权 \(\zeta(z;u,v)\) 的 \((2+8)\) 次显式分解、两条 rigid branches \(\pm\sqrt v\) 与一参数输出位势中的原子 Witt 剥离都已建立之后，还可把输出完全抑制切片与二维/一维原子分量的变量分工进一步压缩为下列两条接口。

\begin{theorem}[二维加权 \texorpdfstring{$\zeta(z;u,v)$}{zeta(z;u,v)} 在 \texorpdfstring{$v=0$}{v=0} 切片上的完全谱塌缩]\label{thm:xi-time-part73a-twoparam-vzero-total-spectral-collapse}
\leanverified{paper\_xi\_time\_part73a\_twoparam\_vzero\_total\_spectral\_collapse}
沿用定理 \ref{thm:real-input-40-det-uv-2plus8} 的记号。则对任意碰撞参数 \(u\) 都有严格恒等式
\[
\boxed{\;
\Delta(z;u,0)=1-z.\;}
\]
因此：
\begin{enumerate}
  \item \(M_{u,0}\) 的非零谱恰为单点 \(\{1\}\)；
  \item 由于 \(M_{u,0}\in M_{40}(\CC)\)，其零特征值的代数重数恰为 \(39\)；
  \item 相应加权 Artin--Mazur \(\zeta\) 函数满足
  \[
  \boxed{\;
  \zeta(z;u,0)=\frac{1}{1-z}.\;}
  \]
  \item 若定义 ghost 坐标与 primitive 生成函数
  \[
  a_n(u,0):=\Tr(M_{u,0}^n),
  \qquad
  P(z;u,0):=\sum_{n\ge 1}p_n(u,0)z^n,
  \]
  则有
  \[
  \boxed{\;
  a_n(u,0)=1\qquad(n\ge 1),\;}
  \]
  以及
  \[
  \boxed{\;
  P(z;u,0)=z,\qquad
  p_1(u,0)=1,\quad p_n(u,0)=0\ (n\ge 2).\;}
  \]
\end{enumerate}
换言之，一旦完全抑制输出位 \(e=1\)，二维模型中的碰撞权 \(u\) 在谱层面完全失明。
\end{theorem}
\begin{proof}
由定理 \ref{thm:real-input-40-det-uv-2plus8}，
\[
\Delta(z;u,v)=(v z^2-1)P_8(z;u,v),
\]
其中
\[
P_8(z;u,v)
=
(u^2 v^4-u^2 v^3)z^8+\cdots +(uv+5v)z^2+z-1.
\]
将 \(v=0\) 代入，则除末两项外其余各项全都消失，从而
\[
P_8(z;u,0)=z-1.
\]
于是
\[
\Delta(z;u,0)=(-1)(z-1)=1-z.
\]
这就给出第一条恒等式。

又由于该加权矩阵作用在真实输入 \(40\) 状态核上，故 \(M_{u,0}\) 为 \(40\times 40\) 矩阵。特征多项式的非平凡部分既然仅为 \(1-z\)，则唯一非零特征值只能是 \(1\)，其余 \(39\) 个特征值全为 \(0\)。于是
\[
\zeta(z;u,0)=\Delta(z;u,0)^{-1}=(1-z)^{-1}.
\]
从而
\[
\log\zeta(z;u,0)=\sum_{n\ge 1}\frac{z^n}{n},
\]
故 ghost 系数满足 \(a_n(u,0)=1\) 对所有 \(n\ge 1\) 成立；同一式也表明 primitive 生成函数恰为
\[
P(z;u,0)=z.
\]
于是 \(p_1(u,0)=1\) 且 \(p_n(u,0)=0\)（\(n\ge 2\)）。证毕。
\end{proof}

\begin{corollary}[二维/一维原子 Witt 校正的变量分裂与单层激活]\label{cor:xi-time-part73a-atomic-witt-bivariate-univariate-variable-rigidity}
分别考虑二维加权模型 \(\zeta(z;u,v)\) 与一参数输出位势模型 \(\zeta(z,u)\)。
\begin{enumerate}
  \item 在二维模型中，若定义 primitive 生成函数
  \[
  P(z;u,v):=\sum_{n\ge 1}p_n(u,v)z^n,
  \qquad
  P_{\mathrm{core}}(z;u,v):=\sum_{n\ge 1}p_n^{\mathrm{core}}(u,v)z^n,
  \]
  以及 ghost 坐标
  \[
  \log\zeta(z;u,v)=\sum_{n\ge 1}\frac{a_n(u,v)}{n}z^n,
  \qquad
  \log\zeta_{\mathrm{core}}(z;u,v)=\sum_{n\ge 1}\frac{a_n^{\mathrm{core}}(u,v)}{n}z^n,
  \]
  则有
  \[
  \boxed{\;
  P(z;u,v)=v z^2+P_{\mathrm{core}}(z;u,v),\;}
  \]
  \[
  \boxed{\;
  a_n(u,v)=
  \begin{cases}
  2v^{n/2}+a_n^{\mathrm{core}}(u,v), & 2\mid n,\\[4pt]
  a_n^{\mathrm{core}}(u,v), & 2\nmid n.
  \end{cases}\;}
  \]
  因而对任意 \(r\ge 1\) 与任意 \(n\ge 1\)，都有
  \[
  \partial_u^r\!\bigl(p_n(u,v)-p_n^{\mathrm{core}}(u,v)\bigr)\equiv 0,
  \qquad
  \partial_u^r\!\bigl(a_n(u,v)-a_n^{\mathrm{core}}(u,v)\bigr)\equiv 0.
  \]
  \item 在一参数输出位势模型中，若记主极点 \(z_\star(u)=\lambda(u)^{-1}\)，则原子 Witt 校正恰在 primitive 的 \(n=2\) 层激活，并满足
  \[
  \boxed{\;
  p_n(u)=p_n^{\mathrm{core}}(u)+u\,\mathbf 1_{\{n=2\}},\;}
  \]
  \[
  \boxed{\;
  a_n(u)=a_n^{\mathrm{core}}(u)+2u^{n/2}\mathbf 1_{\{2\mid n\}},\;}
  \]
  \[
  \boxed{\;
  \log\mathfrak M(u)-\log\mathfrak M_{\mathrm{core}}(u)=u\,z_\star(u)^2.\;}
  \]
  特别地，在无权点 \(u=1\) 上，
  \[
  \log\mathfrak M(1)-\log\mathfrak M_{\mathrm{core}}(1)
  =
  z_\star(1)^2
  =
  \varphi^{-4}.
  \]
\end{enumerate}
\end{corollary}
\begin{proof}
对二维模型，固定任意 \(u\) 并把 \(v\) 视为加权变量。由定理 \ref{thm:real-input-40-det-uv-2plus8}，
\[
\zeta(z;u,v)=\frac{1}{1-vz^2}\,\zeta_{\mathrm{core}}(z;u,v).
\]
于是定理 \ref{thm:xi-finite-euler-primitive-surgery-additivity} 在单原子情形 \((m,E,\ell)=(1,1,2)\) 下直接给出
\[
P(z;u,v)=v z^2+P_{\mathrm{core}}(z;u,v).
\]
另一方面，定理 \ref{thm:xi-time-part73-twoparam-atomic-ghost-u-blindness} 已证明
\[
a_n(u,v)-a_n^{\mathrm{core}}(u,v)=2v^{n/2}\mathbf 1_{\{2\mid n\}}.
\]
由于上述二维原子修正只含 \(v\) 而不含 \(u\)，故关于 \(u\) 的任意高阶导数恒为零。

对一参数模型，全部公式都已由定理 \ref{thm:xi-real-input40-evenlength-abel-shift-rigidity} 给出；最后在 \(u=1\) 处再用该定理中的专门化
\[
z_\star(1)=\varphi^{-2}
\]
便得到
\[
z_\star(1)^2=\varphi^{-4}.
\]
证毕。
\end{proof}

\endinput
